% BHU PYQ -- Manually transcribed & error-corrected
% Paper: BCH-314 : Banking Law & Practice
% Exam: B.Com. (Hons.) Semester V Examination, 2022-23

\documentclass[11pt,a4paper]{article}
\usepackage[a4paper,top=2.5cm,bottom=2.5cm,left=2.8cm,right=2.8cm,headheight=14pt]{geometry}
\usepackage{amsmath,amssymb}
\usepackage[T1]{fontenc}
\usepackage[utf8]{inputenc}
\usepackage{lmodern,microtype,fancyhdr,enumitem,hyperref,lastpage,parskip}

\hypersetup{
    colorlinks=true,
    urlcolor=black,
    linkcolor=black,
    pdftitle={BCH-314: Banking Law \& Practice -- BHU B.Com. (Hons.) Sem V 2022-23}
}

\pagestyle{fancy}
\fancyhf{}
\renewcommand{\headrulewidth}{0.4pt}
\renewcommand{\footrulewidth}{0.4pt}
\fancyhead[L]{\small\textit{Banaras Hindu University}}
\fancyhead[R]{\small\textit{BCH-314: Banking Law \& Practice}}
\fancyfoot[C]{\small Page \thepage\ of \pageref{LastPage}}

\newlist{parts}{enumerate}{2}
\setlist[parts,1]{label=(\alph*),leftmargin=*,itemsep=5pt,topsep=3pt}
\setlist[parts,2]{label=(\roman*),leftmargin=*,itemsep=3pt,topsep=2pt}
\newcommand{\pts}[1]{\hfill\textit{[#1]}}

\begin{document}

\begin{center}
  {\large\bfseries BANARAS HINDU UNIVERSITY}\\[2pt]
  {\normalsize B.Com. (Hons.) --- Semester V Examination, 2022-23}\\[6pt]
  \rule{0.8\linewidth}{0.5pt}\\[6pt]
  {\large\bfseries Paper No. BCH-314: Banking Law \& Practice}\\[4pt]
  {\normalsize Time: 3 Hours\hfill Maximum Marks: 70}
\end{center}

\rule{\linewidth}{0.4pt}

\smallskip
\noindent\textit{Instructions: Attempt all questions. All questions carry equal marks (14 marks each).}

\bigskip

\noindent\textbf{Question 1.}
Define Bank. Explain the working, nature and challenges faced by Commercial Banks in India.\pts{14}

\centerline{\textbf{OR}}

How did Nationalization of Commercial Banks improve bank's efficiency and its significance in economic development? Explain.\pts{14}

\medskip\rule{\linewidth}{0.2pt}

\noindent\textbf{Question 2.}
Discuss the purpose and role of various provisions of the Banking Regulation Act, 1949 in regulating commercial banks in India.\pts{14}

\centerline{\textbf{OR}}

What is the need of RBI? Explain the duties, responsibilities and financial measures adopted by RBI to stabilize the financial system of our country.\pts{14}

\medskip\rule{\linewidth}{0.2pt}

\noindent\textbf{Question 3.}
Point out the origin, objectives, functions and provisions of Regional Rural Banks to promote rural growth in India.\pts{14}

\centerline{\textbf{OR}}

What is EXIM Bank? Discuss the role, functions, features and working of EXIM Bank in India.\pts{14}

\medskip\rule{\linewidth}{0.2pt}

\noindent\textbf{Question 4.}
Define Negotiable Instruments. Discuss its types.\pts{14}

\centerline{\textbf{OR}}

Explain the role, significance and responsibilities of NABARD in agricultural and rural development.\pts{14}

\medskip\rule{\linewidth}{0.2pt}

\noindent\textbf{Question 5.}
Discuss the purpose of different types of loans and advances offered by Commercial Banks to individuals and industries.\pts{14}

\centerline{\textbf{OR}}

Discuss the different types of agricultural finances and its significance in promotion of agricultural sector.\pts{14}

\vfill
\rule{\linewidth}{0.4pt}
\begin{center}\small
\href{https://bhu-exam.vercel.app/}{Click here to visit our website}\\[4pt]
\href{https://me-aryan.vercel.app/}{Click here to visit our contributor}\\[4pt]
\href{https://quashan-me.vercel.app/}{Click here to visit our contributor}
\end{center}

\end{document}
